
\documentclass[11pt,a4paper]{article}
\usepackage[utf8]{inputenc}
\usepackage{float}
\usepackage[T1]{fontenc}
\usepackage[portuguese]{babel}
\usepackage{graphicx}
\usepackage{booktabs}
\usepackage{amsmath,amssymb}
\usepackage{siunitx}
\usepackage{hyperref}
\usepackage{caption}
\usepackage{geometry}
\geometry{top=2cm,bottom=2cm,footskip=1cm}
\usepackage{xcolor}
\usepackage{sectsty}
\usepackage{subcaption}

\definecolor{navyblue}{RGB}{0, 0, 128}
\sectionfont{\color{navyblue}}
\subsectionfont{\color{navyblue}}
\subsubsectionfont{\color{navyblue}}

% --- Capa ---
\begin{document}
\thispagestyle{empty}
\begin{center}
\vspace*{2cm}
\includegraphics[width=0.4\textwidth]{logo_fct.png}\\[1cm]
{\Large \textbf{Simulação de Link de Comunicação em GNU-Radio
Testes de um recetor Low-IF (SDR)}}\\[0.6cm]
{\large CSRF -- 1º Trabalho}\\[2cm]
{\normalsize Diogo Ferreira 73076 \\ Gonçalo Dionísio 73638 \\ Iuri Moças }\\[1cm]
\vspace*{3cm}
{\normalsize Data: 25/05/2026}\\[3.5cm]
{\normalsize Departamento de Engenharia Eletrotécnica e de Computadores}\\
{\normalsize CSRF – 2026/2027 – 1º Semestre}\\
\vfill
\end{center}
\newpage

% --- Sumário ---
\tableofcontents

\listoffigures
\newpage

\section{Introdução}\paragraph{}

Este relatório descreve o trabalho laboratorial desenvolvido no âmbito da unidade curricular de Circuitos 
e Sistemas para Rádio-Frequência, centrado no estudo e teste de sistemas de comunicação por RF. 
O trabalho divide-se em três etapas principais: a simulação a nível de sistema de uma ligação de 
rádio AM com recurso ao GNU Radio; a modelação elétrica e respetiva simulação temporal e espectral 
em SPICE; e o projeto, simulação e caracterização experimental de um filtro passa-banda de 2ª ordem 
para a banda de FM com um Analisador Vetorial de Redes (VNA). Por fim, procedeu-se à integração deste 
filtro num Radio Frequency Front-End (RFFE) para montagem e teste prático de um recetor Low-IF. 
O objetivo central passa por comparar os resultados teóricos, as simulações ideais e com componentes 
reais, e as respetivas medições laboratoriais.

\section{Implementação e Simulação do Link de Comunicação em GNU Radio}\paragraph{}

A primeira parte deste trabalho foi a implementação e simulação de um link de comunicação AM 
utilizando o GNU Radio. O objetivo desta etapa foi compreender o funcionamento de um sistema 
de comunicação AM, desde a modulação até à desmodulação do sinal.



\end{document}
